\chapter{Worked operator products and correlation functions}
\label{ch:correlators}

Characters count states. Correlation functions tell us how those
states interact. The lattice realization is particularly valuable
because it supplies actual fields, positive norms, and exact mode
products. This chapter develops several calculations small enough
to reproduce by hand and explains how symmetry organizes the larger
weight-two three-point tensor.

The correlators below are chiral. Their adjoints are chiral radial
adjoints, not additional antiholomorphic fields. Products with
antiholomorphic counterparts give full vacuum-sector observables.
Generic nonvacuum bulk correlators of the diagonal theory require
the additional sewing data discussed in Chapter~\ref{ch:sewing}.

\section{Normalize before comparing coefficients}
\index{OPE coefficient!normalization}

For orthogonal states $v_i$ with positive norms
$N_i=\langle v_i,v_i\rangle$, write
$\widehat v_i=v_i/\sqrt{N_i}$, taking the positive square root.
For three weight-two primaries, the coefficient of the weight-two
field in their singular OPE is
\begin{equation}
 C_{ij}^{\ k}=
 \frac{\langle v_k,(v_i)_{(1)}v_j\rangle}
      {\sqrt{N_iN_jN_k}}.
 \label{eq:corr-normalization}
\end{equation}
The Hermitian form is antilinear in the first slot. Thus this
coefficient is the correlator with $\widehat v_k^\dagger$ at
infinity. Exact algebraic numerators and positive square roots of
norm products keep both the magnitude and phase of the answer.

The six currents split into affine levels one and thirty-nine.
For their Cartans,
\[
 H_L(z)H_L(0)\sim2z^{-2},\qquad
 H_R(z)H_R(0)\sim78z^{-2},\qquad H_L(z)H_R(0)\sim0.
\]
The defining field $u$ has weight two, charges $(1,13)$, and
transported norm squared twelve. Consequently
$\widehat u=u/\sqrt{12}$ has unit two-point coefficient.
The squared couplings to the unit Cartan currents are
\begin{equation}
 \frac{1^2}{2}=\frac12,
 \qquad \frac{13^2}{78}=\frac{13}{6}.
 \label{eq:corr-current-couplings}
\end{equation}
The unit total stress is $T/2$ at $c=8$, and its coupling to
$\widehat u,\widehat u^\dagger$ is one by the Ward identity.
These numbers provide independent checks on the first terms of
the four-point function.

\section{Fifteen Gaussian terms give a complete four-point function}
\index{four-point function!lattice example}

In the positive unitary exponential basis, the normalized field is
\begin{equation}
 \widehat u=\frac1{\sqrt3}
   \left(e^{a_1}+\zeta e^{a_2}+\bar\zeta e^{a_3}\right),
 \qquad
 \zeta=\frac{\sqrt{13}+3i\sqrt3}{\sqrt{40}},
 \quad |\zeta|=1.
 \label{eq:corr-u-unitary}
\end{equation}
Its three momenta satisfy
\[
 a_i^2=4,\qquad(a_i,a_j)=2\quad(i\ne j).
\]
These inner products, the cocycle convention, and charge
conservation are enough for the calculation. The raw torus-gauge
coefficients in the construction chapter are different because
their exponential basis has squared norms $4,40,40$.

Consider
\[
 G_0(x)=\langle\widehat u^\dagger(\infty)\widehat u(1)
                   \widehat u^\dagger(x)\widehat u(0)\rangle.
\]
Expanding the four sums gives $3^4=81$ assignments. The Gram matrix
of the three momenta has eigenvalues $8,2,2$, so they are linearly
independent. Charge conservation therefore says exactly that
the two positive charges and the two negative charges have the
same multiset. There are three assignments using
one momentum four times and twelve using two distinct momenta.
In the fixed adjoint and cocycle convention, each surviving term
has coefficient $1/9$.

The fifteen terms fall into three groups. The three identical-charge
assignments contribute
$1/[3x^4(1-x)^4]$. The two possible pairing patterns for distinct
charges contribute six terms each, giving
\[
 G_0(x)=\frac1{3x^4(1-x)^4}
       +\frac2{3x^2(1-x)^4}
       +\frac2{3x^4(1-x)^2}.
\]
Putting these over a common denominator yields
\begin{equation}
 G_0(x)=\frac{3-4x+4x^2}{3x^4(1-x)^4}.
 \label{eq:corr-u-neutral}
\end{equation}
This is the full meromorphic chiral correlator, not a truncated
series or a fitted ansatz.

The equal-charge ordering similarly gives
\begin{equation}
 G_+(x)=\langle\widehat u^\dagger(\infty)
    \widehat u^\dagger(1)\widehat u(x)\widehat u(0)\rangle
 =\frac{x^2(4-4x+3x^2)}{3(1-x)^4}.
 \label{eq:corr-u-charged}
\end{equation}
The identities $G_0(1-x)=G_0(x)$ and
$G_+(x)=x^{-4}G_0(1/x)$ are exact crossing checks.
At $x=1/2$, for instance,
$G_0=512/3$ and $G_+=11/3$.

\begin{example}[The first equal-charge exchanged state]
The sum of two distinct momenta has squared norm twelve, so its
exponential has weight six. Equal momenta instead give weight
eight. Hence the leading equal-charge OPE has weight six and
starts at $x^{6-2-2}=x^2$.
The three distinct sums each occur twice in the product of
\eqref{eq:corr-u-unitary}; their coefficients have magnitude
$2/3$. The resulting state has squared norm
$3(2/3)^2=4/3$. It is a weight-six primary, and its normalized
leading coupling can be chosen as $2/\sqrt3$. This reproduces
the leading coefficient of \eqref{eq:corr-u-charged}.
\end{example}

\section{Extracting global blocks by a triangular subtraction}
\index{conformal block!global}

For external weight two, expand the neutral channel in global
conformal blocks:
\begin{equation}
 x^4G_0(x)=A_0+\sum_{n\ge1}A_nx^n
                         {}_2F_1(n,n;2n;x).
 \label{eq:corr-global}
\end{equation}
Each block begins with $x^n$, so its coefficient can be found by
subtracting all lower blocks from the power series. The first
terms of the reduced correlator are
\[
 x^4G_0=1+\frac83x+6x^2+12x^3+\frac{65}3x^4+\cdots.
\]
Thus $A_0=1$, $A_1=8/3$. The weight-one global block is
$x+x^2/2+x^3/3+x^4/4+\cdots$. It contributes $4/3$ at order
$x^2$, leaving $A_2=6-4/3=14/3$. At the next order the known
contributions are $8/9+14/3$, leaving $A_3=58/9$.
Continuing one more step gives
\begin{equation}
 (A_0,A_1,A_2,A_3,A_4)
 =\left(1,\frac83,\frac{14}3,\frac{58}9,\frac{107}{15}\right).
 \label{eq:corr-global-coefficients}
\end{equation}

The first coefficient equals $1/2+13/6$, exactly the two Cartan
exchanges in \eqref{eq:corr-current-couplings}. At weight two,
the total stress contributes one. The other $11/3$ is the sum
of squared couplings to the weight-two Virasoro primaries.
A global coefficient is a sum over orthonormal quasiprimaries;
it is not generally one OPE coefficient or one Virasoro-block
coefficient. This distinction is essential when comparing a
four-point expansion with a table of individual couplings.

\section{Resolving the weight-two space}
\index{weight-two primary tensor}

The $113$ Virasoro primaries consist of $96$ fields in eight
current multiplets and seventeen additional fields. The eight
types are written as dimensions of the two internal left and right
$\mathfrak{sl}_2$ representations. Both current algebras act on
holomorphic fields; these labels do not denote physical chirality.
The types are
\[
 (1,5),(1,9),(1,9),(1,13),(2,2),(2,6),(2,8),(2,14).
\]
Their dimensions sum to
$5+9+9+13+4+12+16+28=96$.
The extra fields comprise an internal right spin-two multiplet $(1,5)$,
a mixed-current multiplet $(3,3)$, and three stress differences.
At internal left affine level one the candidate $(5,1)$ multiplet is null,
so it does not add another five states.

There are four mutually commuting stresses
$(T_L,T_R,T_e,T_f)$ with central charges
\begin{equation}
 \left(1,\frac{117}{41},\frac{2400}{1271},\frac{70}{31}\right).
 \label{eq:corr-four-stresses}
\end{equation}
Their sum is the total stress. A combination $A=\sum_i a_iT_i$
is a total-Virasoro primary precisely when
$\sum_i a_ic_i=0$: this cancels the fourth-order pole in its
OPE with the total stress, while the remaining poles give weight
two and translation. Since the stress space has dimension four,
there are three independent such differences.

\begin{example}[A normalized three-point coefficient from two stresses]
\label{ex:corr-stress-cubic}
Take two commuting stresses and set $A=c_2T_1-c_1T_2$.
Their orthogonal norms give
$N_A=c_1c_2(c_1+c_2)/2$.
The first product is
$A_{(1)}A=2c_2^2T_1+2c_1^2T_2$, so
\[
 \langle A,A_{(1)}A\rangle=2N_A(c_2-c_1).
\]
Consequently
\begin{equation}
 |C_{AAA}|^2
 =\frac{8(c_2-c_1)^2}{c_1c_2(c_1+c_2)}.
 \label{eq:corr-stress-cubic}
\end{equation}
At $c_1=1,c_2=117/41$, this is $23104/9243$.
The calculation checks a nontrivial entry of the large lattice
coupling tensor using only the Virasoro OPE.
\end{example}

Current zero modes act as derivations. If $P$ projects orthogonally
onto the primary space and $F$ is a lowering zero mode, then
\begin{equation}
 P\bigl((Fv_i)_{(1)}v_j\bigr)
 =F P\bigl((v_i)_{(1)}v_j\bigr)
          -P\bigl((v_i)_{(1)}Fv_j\bigr).
 \label{eq:corr-ward-completion}
\end{equation}
The projection commutes with the current action because current
zero modes commute with the total Virasoro algebra and the compact
form acts unitarily. Thus a set of highest-weight seed products
determines all lowered components. In a conventional finite
$\mathfrak{sl}_2$ basis $v_p=F^pv_0$ of dimension $n$,
$Ev_p=p(n-p)v_{p-1}$. Raw torus-gauge norms must still be
transported consistently when using this recursion.

The recorded exact calculation fills the complete $113^3$ tensor,
with $18365$ nonzero ordered entries. The important mathematical
content is the specified primary basis, exact seed products,
equivariant completion, and normalized positive form. A sparse
table is complete only when its unlisted entries have been proved
zero in that defined space. This tensor concerns weight-two
vacuum-module fields; it is not a claim to have computed every
nonvacuum bulk OPE.

\section{Spin-three fields and a null-channel test}
\index{null channel}\index{parafermion algebra!weight-four field}

The Potts and level-thirteen parafermion factors contain a
canonical weight-three field $a$ with intrinsic central charge
$c_0$. To avoid confusion, $c_0$ here denotes a number, not the
central simple-object label used elsewhere. Set
\[
 N=\frac{c_0}{3},\quad b=a_{(1)}a,\quad
 A=31-16\lambda(c_0+2),\quad B=NA.
\]
The parameter pairs are
$(c_0,\lambda)=(4/5,4/13)$ for Potts and
$(8/5,14/473)$ for the parafermion factor.
In the canonical positive normalization, the singular OPE is
\begin{equation}
 a(z)a(0)\sim\frac{N}{z^6}+\frac{2T_0}{z^4}
 +\frac{\partial T_0}{z^3}+\frac b{z^2}
 +\frac{\frac12\partial b-\frac1{12}\partial^3T_0}{z},
 \label{eq:corr-spin3-ope}
\end{equation}
with $a_{(3)}b=Aa$ and $\norm b^2=B$.

One can reconstruct the whole four-point function from these
poles. Define the unnormalized function
$F(x)=\langle a(\infty)a(1)a(x)a(0)\rangle$.
The Ward identity gives
$\langle a(\infty)a(1)T_0(0)\rangle=3N$.
Its first and third derivatives at the last insertion give
$6N$ and $72N$. The corresponding $b$ coefficient and first
derivative are $B$ and $4B$. Therefore the poles at zero have
coefficients $N^2,2c_0,2c_0,B,2B-2c_0$ at orders
$6,4,3,2,1$. Crossing determines the poles at one, and the
identity channel fixes the constant at infinity. Hence
\begin{align}
 F(x)={}&N^2\bigl[1+x^{-6}+(1-x)^{-6}\bigr]\notag\\
 &+2c_0\bigl[x^{-4}+(1-x)^{-4}+x^{-3}+(1-x)^{-3}\bigr]\notag\\
 &+B\bigl[x^{-2}+(1-x)^{-2}\bigr]
 +(2B-2c_0)\bigl[x^{-1}+(1-x)^{-1}\bigr].
 \label{eq:corr-spin3-four}
\end{align}
Subtracting this rational function from the actual correlator
leaves a holomorphic function on the sphere vanishing at infinity,
so the difference is zero. Divide by $N^2$ for unit external fields.

The weight-four quasiprimary descendant of the intrinsic stress is
\[
 \Lambda=:T_0T_0:-\frac3{10}\partial^2T_0,
 \qquad \norm\Lambda^2=\frac{c_0(5c_0+22)}{10}.
\]
Removing it and the derivative from $b$ gives
$W_4=b-\frac{32}{5c_0+22}\Lambda-\frac3{10}\partial^2T_0$.
This subtraction can be checked entirely with Virasoro Gram
products. Put $Q=5c_0+22$. Then
\[
 \norm{\partial^2T_0}^2=20c_0,\qquad
 \langle\Lambda,\partial^2T_0\rangle=0,\qquad
 \norm{\Lambda}^2=\frac{c_0Q}{10}.
\]
Since $W_4$ is orthogonal to these descendants,
\[
 \norm{W_4}^2
 =B-\left(\frac3{10}\right)^2 20c_0
   -\left(\frac{32}{Q}\right)^2\frac{c_0Q}{10}
 =B-\frac{9c_0}{5}-\frac{512c_0}{5Q}.
\]
The squared coupling to the new intrinsic weight-four primary is
$\norm{W_4}^2/N^2$. Substituting $B=(c_0/3)
[31-16\lambda(c_0+2)]$ and simplifying gives
\begin{equation}
 \gamma_4^2=
 \frac{48(c_0+2)\{8-\lambda(5c_0+22)\}}
      {c_0(5c_0+22)}
 =\begin{cases}
 0,&\text{Potts},\\
 60552/2365,&\text{level-thirteen parafermion}.
 \end{cases}
 \label{eq:corr-spin4-test}
\end{equation}
For the parafermion point, $Q=30$ and
$8-\lambda Q=3364/473$, while
$48(c_0+2)/(c_0Q)=18/5$. Their product is the displayed
$60552/2365$. At the Potts point, $Q=26$ and $\lambda Q=8$,
so the same formula vanishes exactly.
In the Potts case this channel is null. In the parafermion case
it is a positive new channel. The test is with respect to the
intrinsic Virasoro algebra of the factor; a field's decomposition
under the total $c=8$ Virasoro algebra is a different question.

\section{A nonvacuum block with fractional powers}
\index{nonvacuum correlator}\index{Gram matrix!light sector}

The lightest $U$-module $c_1$ can be seen by restricting either
fundamental $A_2$ lattice sector. Choose its three weights
\[
 \beta_0=(2/3,1/3),\quad\beta_1=(-1/3,1/3),\quad
 \beta_2=(-1/3,-2/3),
 \qquad(\beta_i,\beta_j)=\delta_{ij}-\frac13.
\]
Each exponential has weight $1/3$. Paired lattice intertwiners
with unit vacuum coefficient give, on $0<x<1$,
\begin{equation}
 \mathcal F_{ij}(x)=
 \langle V_{-\beta_i}(\infty)V_{\beta_i}(1)
              V_{-\beta_j}(x)V_{\beta_j}(0)\rangle
 =x^{-2/3}(1-x)^{-(\beta_i,\beta_j)}.
 \label{eq:corr-light-block}
\end{equation}
The diagonal blocks are $[x(1-x)]^{-2/3}$ and the off-diagonal
ones are $x^{-2/3}(1-x)^{1/3}$. Analytic continuation requires
a path because these are chiral blocks with fractional powers.

The intermediate lattice vacuum restricts to $U\oplus b$.
Here $b$ again denotes the simple module in the modular table,
rather than the preceding weight-four field $a_{(1)}a$.
The internal right Cartan $H_R$ charges are $q_i=(2,0,-2)$.
These form an internal spin-one triplet, whose chiral conformal
weight is $1/3$. The first reduced-block coefficient splits as
\[
 (\beta_i,\beta_j)=\frac{q_iq_j}{78}
       +\left[(\beta_i,\beta_j)-\frac{q_iq_j}{78}\right].
\]
The first term is $U$-current exchange and the second is exchange
through the $b$ top. At $i=j=0$ these are $2/39$ and $8/13$,
whose sum is $2/3$. The matrix of the second terms has
eigenvalues $0,1,35/39$, consistent with positivity.

Restriction of a lattice intertwiner is a genuine $U$ intertwiner.
Thus these are actual nonvacuum chiral data. They are selected
linear combinations of channels determined by the ambient
extension. They do not form a complete channel basis and do not
by themselves determine the diagonal theory's lightest-scalar
bulk four-point function. Multiplying one selected block by its
complex conjugate would be a proposed sewing choice, which still
needs justification.

\section{Exercises}

\begin{exercise}\label{ex:corr-fifteen}
Enumerate the three groups of surviving charge assignments in
the neutral $u$ four-point function and derive their powers of
$x$ and $1-x$ from pairwise momentum inner products.

\par\smallskip\noindent{\footnotesize\normalfont
\hyperref[sol:ex:corr-fifteen]{Solution on page~\pageref*{sol:ex:corr-fifteen}.}}\par
\end{exercise}

\begin{exercise}\label{ex:corr-blocks}
Complete the triangular subtraction leading to $A_4=107/15$.
Use the leading power in the charged channel to determine its
first global-block coefficient.

\par\smallskip\noindent{\footnotesize\normalfont
\hyperref[sol:ex:corr-blocks]{Solution on page~\pageref*{sol:ex:corr-blocks}.}}\par
\end{exercise}

\begin{exercise}\label{ex:corr-stress}
Derive \eqref{eq:corr-stress-cubic} without using a four-point
function. What happens when $c_1=c_2$, and why does that not make
the field $A$ null?

\par\smallskip\noindent{\footnotesize\normalfont
\hyperref[sol:ex:corr-stress]{Solution on page~\pageref*{sol:ex:corr-stress}.}}\par
\end{exercise}

\begin{exercise}\label{ex:corr-spin3}
Verify the Potts zero in \eqref{eq:corr-spin4-test}. Derive the
coefficient of the simple pole in \eqref{eq:corr-spin3-four}
directly from the last term of \eqref{eq:corr-spin3-ope}.

\par\smallskip\noindent{\footnotesize\normalfont
\hyperref[sol:ex:corr-spin3]{Solution on page~\pageref*{sol:ex:corr-spin3}.}}\par
\end{exercise}

\begin{exercise}\label{ex:corr-lightgram}
Write the $3\times3$ matrix
$K_{ij}=\delta_{ij}-1/3-q_iq_j/78$ explicitly. Find its
eigenvectors and eigenvalues without a numerical diagonalization.
Explain why an off-diagonal negative entry is compatible with
positivity of $K$.

\par\smallskip\noindent{\footnotesize\normalfont
\hyperref[sol:ex:corr-lightgram]{Solution on page~\pageref*{sol:ex:corr-lightgram}.}}\par
\end{exercise}
